\documentclass{article}
\usepackage{geometry}
\geometry{a4paper,scale=0.8}
\usepackage[utf8]{inputenc}

\usepackage{amsmath, amssymb}
\usepackage{amsthm}
\usepackage{mathrsfs}
\usepackage{enumitem}
\usepackage{blkarray}
\usepackage{tikz-cd}
\newtheorem{remark}{Remark}
\newtheoremstyle{breakstyle}
    {5pt}
    {10pt}
    {\normalfont}
    {0pt}
    {\bfseries}
    {\newline\indent}
    {0pt}
    {}
\theoremstyle{breakstyle}
\newtheorem{exercise}{Exercise}
\let\ker\relax
\DeclareMathOperator{\ker}{Ker}
\DeclareMathOperator{\coker}{Coker}
\DeclareMathOperator{\coim}{Coim}
\DeclareMathOperator{\im}{Im}
\let\dim\relax
\DeclareMathOperator{\dim}{Dim}
% \let\max\relax
% \DeclareMathOperator{\max}{Max}
% \DeclareMathOperator{\rad}{Rad}
% \DeclareMathOperator{\nil}{Nil}
% \DeclareMathOperator{\obj}{Obj}
% \DeclareMathOperator{\spec}{Spec}
\title{Chapter 2 - Modules}
\author{Flandre}
\date{7 Feb 2026}

\begin{document}
\maketitle

\begin{exercise}
    Since $\gcd(m,n)=1,\exists u,v\in \mathbb{Z},\text{s.t. }$
    $$
    mu+nv=1
    $$

    Now for $\forall x\otimes y\in \left( \mathbb{Z} /m\mathbb{Z} \right) \otimes _{\mathbb{Z}}\left( \mathbb{Z} / n\mathbb{Z} \right) $,
    \begin{align*}
        x\otimes y&= (mu+nv)(x\otimes y) \\
        &= mu(x\otimes y)+nv(x\otimes y) \\
        &= 0 \\
    \end{align*}
    \qed
\end{exercise}

\begin{exercise}
    Consider the exact sequence
    $$
    \mathfrak{a}\xrightarrow{f} A\xrightarrow[]{g} A / \mathfrak{a}\rightarrow 0
    $$
    By \emph{Proposition 2.18}, we have the following exact sequence
    $$
    \mathfrak{a}\otimes M\xrightarrow[]{f\otimes 1} A\otimes M\xrightarrow[]{g\otimes 1}\left( A / \mathfrak{a} \right)  \otimes M\longrightarrow 0
    $$
    Notice there exist trivial isomorphism $\mathfrak{a}\otimes M\cong \mathfrak{a}M,A\otimes M\cong M$.
    So we can have the exact sequence
    $$
    \mathfrak{a}M\xrightarrow[]{\varphi} M\xrightarrow[]{\psi} \left( A / \mathfrak{a} \right) \otimes M\rightarrow 0
    $$

    With $\im{\varphi}=\ker{\psi},\im{\psi}=\left( A / \mathfrak{a} \right) \otimes M$.
    Which is to say
    $$
    \im{\psi}=\left( A / \mathfrak{a} \right) \otimes M,\ker{\psi}=\mathfrak{a}M
    $$
    So by \emph{Module Isomorphism Theorem},
    $$
    \left( A / \mathfrak{a} \right) \otimes M\cong M / \mathfrak{a}M
    $$
    \qed
\end{exercise}

\begin{exercise}
    We follow from the hint:
    For any finitely generated $A$-module $M_0$,
    let $\mathfrak{m}$ be the maximal ideal of $A$ and $k=M_0 / \mathfrak{m}$ be the field derived from quotient.
    Now by \emph{Ex-2.2}, $k\otimes _{A}M_0\cong M_0 / \mathfrak{m}M_0$.
    So by \emph{Nakayama Lemma}, $k\otimes _{A}M_0=0$ will indicate $M_0=0$.

    Back to the problem:
    \begin{align*}
        0&= k\otimes _{A}k\otimes _{A}\otimes _{A}M\otimes _{A}N \\
        &\cong  k\otimes _{A}\left( k\otimes _{A}M \right) \otimes _{A}N \\
        &\cong  k\otimes _{A}M_{k} \otimes _{A}N \\
        &\cong  M_{k}\otimes _{A}\left( k\otimes _{A}N \right)  \\
        &\cong  M_{k}\otimes _{A}N_{k} \\
    \end{align*}
\end{exercise}

\begin{exercise}
    For any $A$-module $N$, consider
    $$
    f:N\longrightarrow N'
    $$
    which is injective.
    Now
    $$
    M\text{ is flat }
    $$
    $\Leftrightarrow$ 
    $$
    f\otimes 1:N\otimes M\longrightarrow N'\otimes M\text{ is injective }
    $$
    $\Leftrightarrow$ (Be considering their kernel)
    $$
    \forall i\in I,f\otimes 1:N\otimes M_{i}\longrightarrow N'\otimes M_{i}\text{ is injective }
    $$
    \qed
\end{exercise}

\begin{exercise}
    We can view $A,xA,x^2A,\ldots $ as a sequence of $A$-module,
    and $A[x]$ is their direct sum where multiplication in the $A$-algebra structure defined in the trivial way.
\end{exercise}

\begin{exercise}
    Define
    $$
    \varphi:A[x]\otimes _{A}M\longrightarrow M[x]
    $$
    by mapping
    $$
    p(x)\otimes m\longmapsto p(x)\cdot m
    $$
    Obviously $\ker{\varphi}=0,\im{\varphi}=M[x]$, so $\varphi$ is an isomorphism.
\end{exercise}

\begin{exercise}
    \leavevmode \vspace{-\baselineskip}
    \begin{enumerate}[label=({\roman*})]
        \item 
            If $a(x),b(x)\in A[x],a(x)b(x)\in \mathfrak{p}[x]$.
            Let $a(x)=a_{n}x^{n}+a_{n-1}x^{n-1}+\ldots +a_1x+a_0,b(x)=b_{n}x^{n}+b_{n-1}x^{n-1}+\ldots +b_1x+b_0$,
            then
            $$
            (a_{n}x^{n}+\ldots +a_1x+a_0)(b_{n}x^{n}+\ldots +b_1x+b_0)\in \mathfrak{p}[x]
            $$
            $$
            a_{n}b_{n}\in \mathfrak{p}\Rightarrow a_{n}\in \mathfrak{p}\text{ or  }b_{n}\in \mathfrak{p}
            $$
            Suppose $a_{n}\in \mathfrak{p}$, then let $a'(x)=a_{n-1}x^{n-1}+a_{n-2}x^{n-2}+\ldots +a_1x+a_0$ and we have
            $$
            a'(x)b(x)\in \mathfrak{p}[x]
            $$
            Repeat the process above and eventually
            $$
            a(x)\in \mathfrak{p}[x]\text{ or  }b(x)\in \mathfrak{p}[x]
            $$
           \qed 
        \item Not necessarily.

            Let $A=k$ be a field and $\mathfrak{m}=(0)$.
            Obviously $\mathfrak{m}[x]=(0)$ is not a maximal ideal in $k[x]$.
    \end{enumerate}
\end{exercise}

\begin{exercise}
    \leavevmode \vspace{-\baselineskip}
    \begin{enumerate}[label=({\roman*})]
        \item 
            Consider the injective $A$-module homomorphism
            $$
            \varphi:R\longrightarrow R'
            $$
            Since $M$ is flat,
            $$
            \varphi\otimes 1:R\otimes _{A}M\longrightarrow R'\otimes _{A}M
            $$
            is injective, therefore
            $$
            \varphi\otimes 1\otimes 1:R\otimes _{A}M\otimes _{A}N\longrightarrow R'\otimes _{A}M\otimes _{A}N
            $$
            is injective, i.e.
            $$
            \varphi\otimes 1:R\otimes _{A}(M\otimes _{A}N)\longrightarrow R'\otimes _{A}(M\otimes _{A}N)
            $$
            is injective.
            So $M\otimes _{A}N$ is flat.
        \item 
            Consider a injective $B$-module homomorphism
            $$
            \varphi:M\longrightarrow M'
            $$
            Since $N$ is a flat $B$-module homomorphism,
            $$
            \varphi\otimes _{B}1:M\otimes _{B}N\longrightarrow M'\otimes _{B}N
            $$
            is injective.
            To prove $N$ is a flat $A$-module, we need to show
            $$
            \varphi\otimes _{A}1:M\otimes _{A}N\longrightarrow M' \otimes _{A}N
            $$
            is injective, i.e. $\ker{\varphi\otimes _{A}1}=\left\{ 0 \right\} $.

            Notice $\ker{\varphi\otimes _{A}1}\subseteq \ker{\varphi\otimes _{B}1}= \left\{ 0 \right\} $, we're done.
            \qed
    \end{enumerate}
\end{exercise}

\begin{exercise}
    Denote
    $$
    0\rightarrow M'\xrightarrow[]{f} M\xrightarrow[]{g} M''\longrightarrow 0
    $$
    We know
    $$
    M\cong \coim{g}\oplus\ker{g}\cong \im{g}\oplus\ker{g}
    $$
    Be the exactness, $\ker{g}=\im{f},\im{g}=M''$, which are both finitely generated.
    So $M$ is finitely generated.
    \qed
\end{exercise}

\begin{exercise}
    Since $M / \mathfrak{a}M\longrightarrow N / \mathfrak{a}N$ is surjective,
    $$
    M\longrightarrow M / \mathfrak{a}M\longrightarrow N / \mathfrak{a}N
    $$
    is surjective.
    Therefore $N=\mathfrak{a}N+u(M)$, and by \emph{Corollary 2.7}, $u(M)=N$.
    \qed
\end{exercise}

\begin{exercise}
    \emph{(The ring multiplication here is Cartesian product!!!)}
    \begin{enumerate}[label=({\roman*})]
        \item When $A^{m}\cong A^{n}$ :
            We use the method in the textbook.

            Let $\mathfrak{m}$ be a maximal ideal of $A$, and denote the isomorphism by
            $$
            \phi:A^{m}\longrightarrow A^{n}
            $$
            By checking the kernel and image, we can easily see that
            $$
            1\otimes _{A}\phi:(A / \mathfrak{m})\otimes _{A}A^{m}\longrightarrow (A / \mathfrak{m})\otimes _{A}A^{n}
            $$
            is an isomorphism by \emph{Proposition 2.18}.
            By \emph{Exercise 2.2} we know $\left( A / \mathfrak{m} \right) ^{m}\cong A^{m} / \mathfrak{m}\cong (A / \mathfrak{m})\otimes A^{m},\left( A / \mathfrak{m} \right) ^{n}\cong A^{n} / \mathfrak{m}\cong (A / \mathfrak{m})\otimes A^{n}$ are two isomorphic linear space over the field $k=A / \mathfrak{m}$.
            Hence
            $$
            \dim{k^{m}}=\dim{k^{n}},m=n
            $$
        \item When $\phi:A^{m}\longrightarrow A^{n}$ is surjective:
            Similar to \emph{(i)}, we obtain a surjective $k$-linear space homomorphism by using \emph{Proposition 2.18}
            $$
            k^{m}\longrightarrow k^{n}
            $$
            So $m>n$ is obtained by the property of linear space homomorphisms.
        \item When $\phi:A^{m}\longrightarrow A^{n}$ is injective:
            Yes.

            Note that we're no longer able to use \emph{Proposition 2.18} to convert ring homomorphisms into linear space homomorphisms.

            Suppose $m>n$, we can define $A^{n}$ to be the submodule of $A^{m}$ by taking the first $n$ terms and extend the range of $\phi$ to $A^{m}$.
            Now
            $$
            \phi(A^{m})\subseteq A^{n}\subseteq A^{m}=\mathfrak{a} A^m
            $$
            where $\mathfrak{a}=A$ is viewed as an ideal.

            Since $A^{m}$ is a finitely generated $A$-module, by \emph{Proposition 2.4}, $\exists a_1,a_2,\ldots ,a_{s}\in A,\text{s.t. }$
            $$
            \phi^{s}+a_{1}\phi^{s-1}+\ldots +a_{s}=0,a_{s}\neq 0
            $$
            Specifically,
            $$
            (\phi^{s}+a_1\phi^{s-1}+\ldots +a_{s-1}\phi+a_{s}\mathrm{Id})(0,\ldots ,0,1)=0
            $$
            
            Notice $\forall a\in A^{m}$, because $m>n$, the last term of $\phi(a)$ must be 0.
            So
            $$
            (\phi^{s}+\ldots +a_{s-1}\phi)(0,\ldots ,0,1)=(\ldots ,0)
            $$
            But
            $$
            a_{s}\mathrm{Id}(0,\ldots ,0,1)=(\ldots ,a_{s})
            $$
            Hence, it's impossible to have $(\phi^{s}+a_1\phi^{s-1}+\ldots +a_{s-1}\phi+a_{s}\mathrm{Id})(0,\ldots ,0,1)=0$
            So $m\le n$.
    \end{enumerate}
    \qed
\end{exercise}

\begin{exercise}
    Let $A^{n}=\left<e_1,e_2,\ldots ,e_{n} \right> $ where
    $$
    \begin{cases}
        e_1=\left( 1,0,0,\ldots ,0 \right) \\
        e_2=\left( 0,1,0,\ldots ,0\right)\\
        \vdots
        e_{n}=\left( 0,0,0,\ldots ,1 \right) 
    \end{cases}
    $$
    and choose $u_{k}\in M,\text{s.t. }\phi(u_{k})=e_{k}$.
    We can prove every element in $M$ can be expressed in the sum of element in $\left<u_1,u_2,\ldots ,u_{n} \right> ,\ker{\phi}$.
    Also, $\ker{\phi}\cap \left<u_1,\ldots,u_{n} \right> =\left\{ 0 \right\} $,
    so $M=\ker{\phi}\oplus \left<u_1,\ldots ,u_{n} \right>$.
\end{exercise}

\begin{exercise}
    \leavevmode \vspace{-\baselineskip}
    \begin{enumerate}[label=({\roman*})]
        \item $\ker{g}=0$, so $g$ is injective.
        \item Define
            $$
            p:N_{B}\longrightarrow N
            $$
            by
            $$
            b\otimes y\longmapsto by
            $$
            and it's trivial that
            $$
            N_{B}=\im{g}\oplus \ker{p}
            $$
    \end{enumerate}
    \qed
\end{exercise}

\begin{exercise}
    To prove $\mu_{i}=\mu_{j}\circ \mu_{ij}$,
    we only need to understand that in the structure of $M=C / D$,
    elements are the equivalence classs in $C$ where equivalence relation defined by homomorphisms.
    The action of $\mu_{ij}$ is simply sending one element to another in the same equivalence classs.

    After one understand this, the problem becomes rather simple.
\end{exercise}

\begin{exercise}
    \leavevmode \vspace{-\baselineskip}
    \begin{enumerate}[label=({\roman*})]
        \item This is to say that $\mu:C\longrightarrow C / D$ is surjective, which is obviously correct.
        \item This is automatically obtained by the definition of $D$.
    \end{enumerate}
\end{exercise}

\begin{exercise}
    Because any such homomorphism that satisfy
    $$
    \forall i\in I,\alpha_{i}=\alpha\circ \mu_{i}
    $$
    is unique determined by the inner structure of $M_{i},N,M$.
\end{exercise}

\begin{exercise}
    Because in the notation of \emph{Exercise 2.14},
    $D=\left\{ 0 \right\} $ and therefore
    $$
    \lim_{\rightarrow }M_{i}=C / D=C=\sum_{i\in I}^{} M_{i} 
    $$
    \qed
\end{exercise}

\begin{exercise}
    $$
        \begin{tikzcd}
        M_i \arrow[r, "\phi_i"] \arrow[d, "\mu_i"] & N_i \arrow[d, "\nu_i"] \\
        M \arrow[r, "\phi"]                        & N                     
        \end{tikzcd}
    $$
    Because form those $\nu_{i}\circ \phi_{i}$ defining the direct limit homomorphism $M_{i}\longrightarrow N$,
    the unique homomorphism $\phi$ is obtained by the univeral property of direct limit.
\end{exercise}

\begin{exercise}
    $$
        \begin{tikzcd}
        M_i \arrow[r, "\phi_i"] \arrow[d, "\mu_i"] & N_i \arrow[d, "\nu_i"] \arrow[r, "\psi_i"] & P_i \arrow[d, "\xi_i"] \\
        M \arrow[r, "\phi"]                        & N \arrow[r, "\psi"]                        & P                     
        \end{tikzcd}
    $$
    By \emph{Exercise 2.15}, each elements in $M$ can be expressed in the form $\mu_{i}(x_{i})$,
    and $\ker{\psi}=\left\{ \nu_{i} \ker{\psi_{i}}\right\}_{i\in I} $.
    Therefore
    \begin{align*}
        \im{\phi}&= \left\{ \im{\phi\circ\mu_{i} } \right\}_{i\in I}  \\
                 &= \left\{ \im{\nu_{i}\circ \phi_{i}} \right\}_{i\in I}  \\
                 &= \left\{ \nu_{i}\im{\phi_{i}} \right\}_{i\in I}  \\
                 &= \left\{ \nu_{i}\ker{\psi_{i}} \right\}_{i\in I}  \\
                 &= \ker{\psi} \\
    \end{align*}
    \qed
\end{exercise}

\begin{exercise}
    $$
        \begin{tikzcd}
        M_i\times N \arrow[r] \arrow[d]            & M_i\otimes N \arrow[d]            \\
        (\lim_{\rightarrow} M_i)\times N \arrow[r] & \lim_{\rightarrow}(M_i \otimes N)
        \end{tikzcd}
    $$
    
\end{exercise}

\end{document}
